\documentclass[11pt,a4paper]{article}
\usepackage[utf8]{inputenc}
\usepackage[T1]{fontenc}
\usepackage{lmodern}
\usepackage[french]{babel}
\usepackage{amsmath,amssymb,mathtools}
\usepackage{icomma}
\usepackage{xcolor}
\usepackage{colortbl}
\usepackage{tikz}
\usepackage[most]{tcolorbox}
\usepackage{enumitem}
\usepackage{array}
\usepackage{booktabs}
\usepackage{fancyhdr}
\usepackage[hidelinks]{hyperref}
\usetikzlibrary{arrows.meta,positioning,calc,patterns,backgrounds,shapes.geometric,decorations.pathreplacing,decorations.markings}
\usepackage[a4paper,margin=2.15cm,headheight=26pt,headsep=12pt]{geometry}
\usepackage{titlesec}

% ---------------- Palette ----------------
\definecolor{primary}{RGB}{0,151,169}
\definecolor{primarydark}{RGB}{0,88,101}
\definecolor{defcol}{RGB}{0,114,178}
\definecolor{thmcol}{RGB}{0,140,120}
\definecolor{formulacol}{RGB}{198,93,9}
\definecolor{excol}{RGB}{0,135,90}
\definecolor{attncol}{RGB}{200,40,55}
\definecolor{methcol}{RGB}{90,90,100}
\definecolor{remarkcol}{RGB}{120,94,200}
\definecolor{barcol}{RGB}{0,151,169}
\definecolor{barcold}{RGB}{0,110,124}
\definecolor{modecol}{RGB}{200,55,55}
\definecolor{meancol}{RGB}{0,90,200}
\definecolor{medcol}{RGB}{160,90,190}
\definecolor{quartcol}{RGB}{0,135,90}
\definecolor{ogivecol}{RGB}{176,74,0}

% ---------------- Boîtes ----------------
\tcbset{boxbase/.style={enhanced,breakable,boxrule=0.9pt,arc=2.5pt,
  left=3mm,right=3mm,top=2.3mm,bottom=2.3mm,
  fonttitle=\bfseries,coltitle=white,toptitle=0.6mm,bottomtitle=0.6mm}}

\newtcolorbox{definition}[1][]{boxbase,colback=defcol!6,colframe=defcol,
  title={\textsc{Définition}\ifx\relax#1\relax\relax\else\textmd{ : \itshape #1}\fi}}
\newtcolorbox{propriete}[1][]{boxbase,colback=thmcol!7,colframe=thmcol,
  title={\textsc{Propriété}\ifx\relax#1\relax\relax\else\textmd{ : \itshape #1}\fi}}
\newtcolorbox{formule}[1][]{boxbase,colback=formulacol!7,colframe=formulacol,
  title={\textsc{Formule}\ifx\relax#1\relax\relax\else\textmd{ : \itshape #1}\fi}}
\newtcolorbox{exemple}[1][]{boxbase,colback=excol!6,colframe=excol,
  title={\textsc{Exemple}\ifx\relax#1\relax\relax\else\textmd{ : \itshape #1}\fi}}
\newtcolorbox{methode}[1][]{boxbase,colback=methcol!8,colframe=methcol,sharp corners,
  title={\textsc{Méthode}\ifx\relax#1\relax\relax\else\textmd{ : \itshape #1}\fi}}
\newtcolorbox{remarque}[1][]{boxbase,colback=remarkcol!7,colframe=remarkcol,
  title={\textsc{Astuce}\ifx\relax#1\relax\relax\else\textmd{ : \itshape #1}\fi}}
\newtcolorbox{attention}[1][]{boxbase,colback=attncol!7,colframe=attncol,
  title={\textsc{Attention}\ifx\relax#1\relax\relax\else\textmd{ : \itshape #1}\fi}}

% Boîte "exercice" et "corrigé" (titre libre passé en argument)
\newtcolorbox{exobox}[1]{boxbase,colback=primary!4,colframe=primary,
  before skip=8pt,after skip=8pt,title={#1}}
\newtcolorbox{corbox}[1]{boxbase,colback=excol!5,colframe=excol!85!black,
  before skip=8pt,after skip=8pt,title={#1}}

% ---------------- En-tête ----------------
\newcommand{\docpart}[1]{\def\thedocpart{#1}}
\docpart{}
\pagestyle{fancy}
\fancyhf{}
\fancyhead[L]{\small\textcolor{primarydark}{\textbf{Fiche 8} — Statistiques descriptives}}
\fancyhead[R]{\small\textcolor{primarydark}{\textsc{\thedocpart}}}
\fancyfoot[C]{\small\thepage}
\renewcommand{\headrulewidth}{0.8pt}
\renewcommand{\headrule}{\hbox to\headwidth{\color{primary}\leaders\hrule height \headrulewidth\hfill}}

\titleformat{\section}{\large\bfseries\color{primarydark}}{\thesection}{0.6em}{}
\titleformat{\subsection}{\normalsize\bfseries\color{primary}}{\thesubsection}{0.6em}{}
\setlist{topsep=2pt,itemsep=1.5pt,parsep=0pt}

% Bandeau de titre du document
\newcommand{\bandeau}[2]{%
\begin{tcolorbox}[boxbase,colback=primary,colframe=primarydark,halign=center,
   before skip=2pt,after skip=12pt]
{\LARGE\bfseries\color{white} #1}\\[3pt]{\normalsize\color{white} #2}
\end{tcolorbox}}

\newcommand{\ecm}{\sigma_m}
\newcommand{\ect}{\sigma}
\newcommand{\med}{M\!e}
\docpart{Tutoriel — Thème 4}
\begin{document}
\bandeau{Thème 4 — Les indicateurs de dispersion}{Étendue, écart moyen, variance, écart type}

\noindent\textbf{But du thème.} Deux séries de \emph{même moyenne} peuvent être très différentes :
$\{10,10,10\}$ et $\{0,10,20\}$ ont la même moyenne $10$, mais l'une est concentrée, l'autre étalée.
Les indicateurs de \textbf{dispersion} mesurent cet étalement autour de la moyenne. Ce thème
automatise les quatre : étendue, écart moyen, variance (deux formules) et écart type.

\section{L'étendue : la mesure la plus simple}

\begin{formule}[étendue]
\[ \text{étendue}=x_{\max}-x_{\min}. \]
La plus grande valeur moins la plus petite. Simple, mais \emph{grossière} : elle ne regarde que les
deux extrêmes.
\end{formule}

\section{L'écart moyen : la distance moyenne à la moyenne}

\begin{formule}[écart moyen $\ecm$]
\[
\boxed{\;\ecm=\frac{1}{N}\sum_{i=1}^{k} n_i\,|m-x_i|\;}
\]
On prend la \textbf{valeur absolue} de chaque écart $m-x_i$ car, sans elle, la somme pondérée des
écarts vaut toujours $0$ (positifs et négatifs se compensent). Unité : celle du caractère.
\end{formule}

\begin{methode}[calculer $\ecm$ en colonnes]
Colonne $|m-x_i|$, puis colonne $n_i|m-x_i|$, on additionne, on divise par $N$.
\end{methode}

\section{La variance et l'écart type}

Pour une mesure « algébriquement meilleure », on élève les écarts \textbf{au carré}.

\begin{formule}[variance --- deux formules équivalentes]
\[
\text{(directe)}\quad V=\frac{1}{N}\sum n_i\,(m-x_i)^2
\qquad\Longleftrightarrow\qquad
\text{(König--Huygens)}\quad \boxed{\,V=\frac{1}{N}\sum n_i\,x_i^{\,2}-m^2\,}
\]
« La moyenne des carrés moins le carré de la moyenne. » Unité : le \textbf{carré} de celle du
caractère. Toujours $V\geqslant0$.
\end{formule}

\begin{formule}[écart type $\ect$]
\[ \boxed{\;\ect=\sqrt{V}\;} \]
Racine de la variance : on revient à l'unité du caractère. C'est l'indicateur de dispersion
le plus utilisé.
\end{formule}

\begin{remarque}[pourquoi préférer König--Huygens]
La formule directe fait intervenir $m$ (souvent arrondi) sur \emph{chaque} ligne : l'erreur se
propage. König--Huygens n'utilise $m$ qu'une fois (le $-m^2$ final) : moins d'arrondis, une seule
colonne $n_i x_i^2$ à remplir. Au concours, c'est presque toujours la plus rapide.
\end{remarque}

\begin{exemple}[mini-calcul complet]
Série discrète $x=2\,(n{=}1),\ 4\,(n{=}2),\ 6\,(n{=}1)$, $N=4$.
Moyenne : $m=\dfrac{2+8+6}{4}=4$.
König--Huygens : $\sum n_i x_i^2=1(4)+2(16)+1(36)=72$, donc
\[ V=\frac{72}{4}-4^2=18-16=2,\qquad \ect=\sqrt2\approx1{,}41. \]
\end{exemple}

\begin{attention}[deux contrôles de cohérence]
\begin{itemize}
  \item $V\geqslant0$ \emph{toujours} : une variance négative signale une erreur de calcul.
  \item $\ect\geqslant\ecm$ \emph{toujours} : l'écart type est supérieur (ou égal) à l'écart moyen
        (le carré « surpondère » les grands écarts).
  \item Ne pas oublier : $V$ est au \textbf{carré} de l'unité ; $\ect$ et $\ecm$ dans l'unité simple.
\end{itemize}
\end{attention}
\end{document}
